\section{Introduction}

\subsection{Purpose}

This specification states what \texttt{fsb} does, precisely enough that a person
could write an equivalent program from it and a tester could check the program
against it.  A statement belongs here if a user, or a tester who cannot see the
source, could confirm it.  How the behavior is achieved, and why that way was
chosen, belongs to the design specification.

\subsection{Product summary}

\texttt{fsb} serves a web interface for browsing the user's own files on macOS.
It listens on the loopback address only, opens the interface in a browser, and
lets the user list folders, filter and search by name, and preview files.  It
never writes to, renames, moves or deletes anything it serves.  Two rule files let
the user hide clutter (\emph{ignore}) and permanently block sensitive locations
(\emph{deny}); a built-in set of deny rules for known credential locations is
active even when no file exists.

\subsection{Goals}

\begin{description}
  \item[G1] Fast, comfortable browsing of an entire home folder in a browser,
        including folders with 100{,}000 or more entries.
  \item[G2] \emph{Local only}: no exposure beyond the loopback interface, with
        defenses against attacks that a web page can mount on a local server.
  \item[G3] \emph{Read-only}: nothing that \texttt{fsb} serves can be altered
        through it.
  \item[G4] Explicit, testable exclusion semantics: a two-tier hide and deny model.
  \item[G5] Simple distribution: a single binary.
  \item[G6] Growth in stages (crawl, walk, run) without redesigning the security
        core.
\end{description}

\subsection{Non-goals}

\begin{itemize}
  \item Remote access, sharing, multiple users, or any authentication beyond the
        per-launch credential of Section~\ref{sec:security}.
  \item Any operation that changes files: upload, rename, move, delete, edit.
  \item Platforms other than macOS in this version.
  \item Replacing Finder or a terminal as a file manager.
  \item Running as a background service.
\end{itemize}

\subsection{Users}

The primary user is the author, on their own Mac.  Secondary users are developers
and power users who install the program from a public repository and need to trust
it after reading its documentation: the manual page, the README, and the security
policy.

\subsection{Terms}

\begin{description}
  \item[Root] A folder that \texttt{fsb} serves.  Nothing outside the roots is
        reachable.
  \item[Entry] A file, folder or link inside a folder.
  \item[Ignore (hide)] A rule that removes an entry from listings and search
        results.  The entry stays reachable by its path.
  \item[Deny (block)] A rule that makes a path unreachable by any means.  A denied
        path is answered exactly as a path that does not exist.
  \item[Core rules] The built-in deny rules for known credential and secret
        locations.
  \item[Launch URL] The single-use address that \texttt{fsb} prints at startup.
  \item[Session] The browser's authenticated state after the launch URL has been
        used once.
  \item[Cloud-only file] A file whose contents are stored only in a cloud service
        and not on this machine (an ``optimized storage'' placeholder).
\end{description}

\subsection{Conventions}

\begin{itemize}
  \item A requirement is written \emph{shall}.  Each has an identifier such as
        \rid{CLI-4}, which is a link in the PDF; the identifiers are stable and are
        cited by tests and by the other specifications.  An identifier is never
        reused.
  \item ``The user'' is the person at the keyboard; ``the browser'' is the
        browser they opened the interface in.
  \item Paths are absolute unless stated.  ``Case'' and ``the same name''
        are defined in Section~\ref{sec:names}.
  \item Limits are collected in Section~\ref{sec:limits}.
\end{itemize}

\subsection{Related documents}

\begin{description}
  \item[Design specification] How the behavior is achieved, and why.
  \item[Algebraic specification] The laws that the access rules and the pure
        functions obey, each tied to an executable check.
  \item[Security policy] The threat model, the guarantees, the known limitations,
        and how they were tested.
  \item[Manual page] \code{fsb(1)}: the command line, files and exit status, in the
        form users expect.
\end{description}
